\documentclass[conference]{IEEEtran}
\IEEEoverridecommandlockouts
% The preceding line is only needed to identify funding in the first footnote. If that is unneeded, please comment it out.

\usepackage{cite}
\usepackage{amsmath,amssymb,amsfonts}
\usepackage{algorithmic}
\usepackage{graphicx}
\usepackage{textcomp}
\usepackage{xcolor}
\usepackage{url}
\def\BibTeX{{\rm B\kern-.05em{\sc i\kern-.025em b}\kern-.08em
    T\kern-.1667em\lower.7ex\hbox{E}\kern-.125emX}}

\begin{document}

\title{Data Center Energy Supply Modeling Using Historical\\
Price and Demand Data with Microgrid Integration%
\thanks{Course project, ELEN 4510: Modern \& Clean Energy Systems,
Columbia University, Spring 2026.}}

\author{
\IEEEauthorblockN{Afroditi Fragkiadaki, Daniel Holland, Raphael Vogeley, Tselmeg Mendsaikhan}
\IEEEauthorblockA{\textit{Department of Electrical Engineering, Columbia University, New York, NY, USA}\\
af3619@columbia.edu,\ doh2105@columbia.edu,\ rpv2113@columbia.edu,\ tm3516@columbia.edu}
}

\maketitle

% -----------------------------------------------------------------------
\begin{abstract}
We present an open-source hourly simulation and optimization framework for
data center energy supply planning.
The user specifies two key parameters: the facility IT load (MW) and a
minimum on-site renewable energy fraction (\%).
Given those inputs, the tool evaluates and compares two supply architectures
across any combination of U.S.\ ISO/RTO markets: an islanded behind-the-meter
microgrid composed of solar PV, battery storage, and natural gas reciprocating
engines, and a grid-connected architecture that supplements on-site resources
with wholesale electricity imports.
For each market, the framework constructs weather-driven load, solar, battery,
gas dispatch, and day-ahead price time series, then searches over candidate
asset sizes to find the minimum system levelized cost of energy (sLCOE) design
that meets the renewable constraint.
We present results for a representative 50\,MW data center facility with a
30\% renewable floor across all seven ISO/RTO markets (ERCOT, PJM, MISO,
CAISO, SPP, NYISO, and ISO New England), showing that the grid-connected
architecture is cheaper in every market studied.
We then perform a detailed analysis for ERCOT (Dallas-Fort Worth), examining
the full sLCOE surface, the cost premium imposed by the renewable constraint,
the sensitivity of results to key techno-economic parameters, and the
conditions under which islanded operation could become cost-competitive.
\end{abstract}

\begin{IEEEkeywords}
data centers, microgrids, solar PV, battery energy storage, natural gas
generation, system LCOE, grid optimization, renewable energy constraint
\end{IEEEkeywords}

% -----------------------------------------------------------------------
\section{Introduction}
\label{sec:intro}

Data centers have become material electricity-sector loads.
The International Energy Agency estimates that data centers consumed
approximately 415\,TWh globally in 2024---roughly 1.5\% of global electricity
consumption---and projects this figure could approximately double by 2030
under its base case \cite{b1}.
This growth is driven by cloud services, accelerated computing, and
artificial-intelligence workloads.
Data centers are also geographically concentrated: even when individual global
shares appear modest, local transmission-queue congestion and wholesale-price
impacts can be significant.

This context motivates a practical planning question for large computing
loads: should a data center rely only on the bulk grid, build a
behind-the-meter (BTM) microgrid, or combine both?
Microgrids are attractive because they can improve resilience, integrate local
renewable generation, and support islanded operation during grid
disturbances \cite{b2}.
At the same time, a microgrid that must serve a data center without grid
support in every hour must be sized for reliability, which forces overbuilding
of dispatchable generation and storage relative to normal operating needs.

The project modeled this trade-off by comparing two strategies.
\emph{Model A} is an islanded BTM microgrid with solar PV, battery energy
storage systems (BESS), and natural gas reciprocating internal combustion
engines (RICE).
\emph{Model B} is grid-connected and allows hourly imports from regional
electricity markets.
The comparison spans ERCOT, PJM, MISO, CAISO, SPP, NYISO, and ISO New England.

Our contribution is an open-source decision tool that connects public weather,
solar, market-price, and technology-cost data to a technology-cost sLCOE
optimizer.\footnote{Live tool:
\url{https://afroditifragiadaki-datacenter-microgrid-sim-dashboard-zfcrhl.streamlit.app/}}
The tool is not intended to replace production-grade interconnection studies
or stochastic resource adequacy analysis.
Instead, it provides a transparent first-order comparison of energy
architectures under a consistent set of assumptions.

% -----------------------------------------------------------------------
\section{Project Scope and System Architectures}
\label{sec:scope}

The analysis assumes an existing data center with configurable IT load.
The baseline run uses 50\,MW of IT load; facility load is larger because
cooling and auxiliary systems are captured through Power Usage Effectiveness
(PUE).
All technologies are modeled behind the meter except utility imports in the
grid-connected case.

\subsection{Islanded Microgrid (Model A)}

The islanded configuration has no grid connection.
Solar generation serves load first; surplus solar charges the BESS; any
remaining deficit is served by natural gas.
Because the data center must be served in every modeled hour, gas capacity is
the reliability backstop sized to eliminate unserved energy over the full
8,784-hour simulation.

\subsection{Grid-Connected (Model B)}

The grid-connected configuration uses the same solar-first dispatch, but any
residual demand after solar and BESS is covered by the cheaper of on-site gas
and grid imports each hour, using the day-ahead locational marginal price
(LMP).
The grid is treated as an infinite-capacity source; unserved energy is always
zero in Model B.

% -----------------------------------------------------------------------
\section{Data Sources and Preprocessing}
\label{sec:data}

\textbf{Weather.}
Hourly dry-bulb temperature is retrieved from the Open-Meteo Historical
Weather API \cite{b5} for a representative datacenter location in each ISO
(Table~\ref{tab:iso}).
Data covers all 8,784 hours of the 2024 leap year.

\textbf{Solar resource.}
Solar irradiance and AC output are obtained from the NREL PVWatts V8
API \cite{b3}, which uses NSRDB TMY data internally.
The TMY 8,760-hour profile is expanded to 8,784 hours by duplicating
February 28 as February 29.
Site-specific tilt angles (set to latitude) and south-facing azimuth
(180$^\circ$) are used for all ISOs.
Module type is ``Premium'' ($\gamma = -0.35$\%/$^\circ$C), DC/AC ratio 1.2,
system losses 15\%.

\textbf{Wholesale electricity prices.}
Hourly real-time LMPs for 2024 are retrieved from each ISO's published data
portal.
Negative prices are retained for dispatch decisions but clipped to zero in
cost accounting.

\textbf{Technology costs.}
Capital and operating costs are sourced from Lazard LCOE+ v18.0 \cite{b6},
NREL Annual Technology Baseline 2025 \cite{b7}, and the EIA Short-Term
Energy Outlook \cite{b8}.
Gas fuel prices use 2024 EIA industrial state-level prices with estimated
basis differentials (Table~\ref{tab:iso}).

\begin{table}[htbp]
\caption{ISO Locations and Regional Parameters}
\label{tab:iso}
\begin{center}
\begin{tabular}{|c|c|c|c|}
\hline
\textbf{ISO} & \textbf{Location} & \textbf{\begin{tabular}[c]{@{}c@{}}Gas\\(\$/MMBtu)\end{tabular}} & \textbf{\begin{tabular}[c]{@{}c@{}}CAPEX\\mult.\end{tabular}} \\
\hline
ERCOT  & Dallas-Fort Worth, TX & 2.50 & 1.00 \\
\hline
PJM    & Pittsburgh, PA        & 3.20 & 1.05 \\
\hline
MISO   & Indianapolis, IN      & 3.00 & 1.03 \\
\hline
CAISO  & Fresno, CA            & 5.00 & 1.15 \\
\hline
SPP    & Oklahoma City, OK     & 2.80 & 1.00 \\
\hline
NYISO  & Albany, NY            & 5.50 & 1.18 \\
\hline
ISO-NE & Hartford, CT          & 6.00 & 1.20 \\
\hline
\end{tabular}
\end{center}
\end{table}

\begin{table}[htbp]
\caption{Model Input Parameters}
\label{tab:inputs}
\begin{center}
\begin{tabular}{|l|c|c|}
\hline
\textbf{Parameter} & \textbf{Value} & \textbf{Source} \\
\hline
IT load baseline & 50 MW & --- \\
Annual demand (ERCOT) & 522,524 MWh & Simulated \\
\hline
\multicolumn{3}{|l|}{\textit{Solar PV (utility-scale, fixed tilt)}} \\
\hline
CAPEX & \$950/kW$_{\rm DC}$ & NREL ATB 2025 \\
Fixed O\&M & \$10/kW/yr & NREL ATB 2025 \\
Lifetime & 30 yr & NREL ATB 2025 \\
\hline
\multicolumn{3}{|l|}{\textit{Battery Energy Storage (4-hr)}} \\
\hline
CAPEX & \$280/kWh & Lazard LCOS V18 \\
Charge / discharge eff. & 96\% each & Literature \\
Min.\ SoC & 20\% & Manufacturer \\
Augmentation & 2.5\%/yr & Lazard LCOS V18 \\
Lifetime & 15 yr & Lazard LCOS V18 \\
\hline
\multicolumn{3}{|l|}{\textit{Natural Gas RICE}} \\
\hline
CAPEX & \$1,100/kW & NREL ATB 2025 \\
Fixed O\&M & \$15/kW/yr & NREL ATB 2025 \\
Variable O\&M & \$5/MWh & NREL ATB 2025 \\
Heat rate & 9.0 MMBtu/MWh & Industry \\
CO$_2$ factor & 0.0531 t/MMBtu & EPA AP-42 \\
Lifetime & 20 yr & NREL ATB 2025 \\
\hline
\multicolumn{3}{|l|}{\textit{Financial}} \\
\hline
Discount rate $r$ & 8\% & Project finance \\
Project life $n$ & 25 yr & Project finance \\
\hline
\end{tabular}
\end{center}
\end{table}

% -----------------------------------------------------------------------
\section{Mathematical Modeling}
\label{sec:model}

\subsection{Demand Model}

Total hourly facility load $L_t$ (MW) is the product of constant IT load
$L_{\rm IT} = 50$\,MW and the Power Usage Effectiveness at ambient
temperature $T_t$ ($^\circ$C):
%
\begin{equation}
  L_t = L_{\rm IT} \times \mathrm{PUE}(T_t).
  \label{eq:load}
\end{equation}

A naive step-function PUE creates discontinuous load jumps at threshold
crossings that are physically unrealistic.
A piecewise-linear interpolation over five control points is used instead:
%
\begin{equation}
  \mathrm{PUE}(T) =
  \begin{cases}
    1.12 & T \leq 10\,^\circ\mathrm{C} \\[2pt]
    1.12 + \dfrac{T-10}{1250} & 10 < T \leq 25\,^\circ\mathrm{C} \\[6pt]
    1.20 + \dfrac{T-25}{833} & 25 < T \leq 35\,^\circ\mathrm{C} \\[6pt]
    1.32 + \dfrac{T-35}{769} & T > 35\,^\circ\mathrm{C}
  \end{cases}
  \label{eq:pue}
\end{equation}

\noindent clamped at $\mathrm{PUE}_{\min}=1.12$ (free economizer cooling,
$T\leq 10\,^\circ$C) and $\mathrm{PUE}_{\max}=1.45$ (full mechanical
chilling, $T\geq 60\,^\circ$C).
For ERCOT (Dallas-Fort Worth) the simulated average PUE is 1.190, average
facility load 59.5\,MW, and annual consumption 522,524\,MWh.

\subsection{Solar Generation Model}

Hourly AC solar output (MW) is
%
\begin{equation}
  P_{\rm sol,t} = S \cdot \mathrm{CF}_t,
  \label{eq:solar}
\end{equation}

\noindent where $S$ (MW$_{\rm DC}$) is the installed DC nameplate (optimizer
decision variable) and $\mathrm{CF}_t \in [0,1]$ is the hourly capacity
factor from PVWatts.
For ERCOT the annual average CF is 17.05\% with an AC yield of
1,492\,kWh/kW.

\subsection{Battery Energy Storage Model}

Battery state-of-charge $E_t$ (MWh) evolves as
%
\begin{equation}
  E_t = E_{t-1}
      + P_{\rm ch,t}\,\eta_{\rm ch}\,\eta_{{\rm temp},t}
      - \frac{P_{\rm dis,t}}{\eta_{{\rm temp},t}\,\eta_{\rm dis}},
  \label{eq:bess}
\end{equation}

\noindent subject to
%
\begin{align}
  \mathrm{SoC}_{\min}\,B &\leq E_t \leq B, \label{eq:soc} \\
  0 \leq P_{\rm ch,t}\,,\; P_{\rm dis,t} &\leq B/\tau, \label{eq:power}
\end{align}

\noindent where $B$ (MWh) is battery energy capacity (optimizer decision
variable), $\tau=4$\,h is the energy-to-power ratio,
$\eta_{\rm ch} = \eta_{\rm dis} = 0.96$, and
$\mathrm{SoC}_{\min} = 0.20$ (80\% depth of discharge).
The nominal round-trip efficiency is
%
\begin{equation}
  \eta_{\rm rt} = \eta_{\rm ch}\,\eta_{\rm dis}
               = 0.96 \times 0.96 = 92.2\%.
  \label{eq:eta_rt}
\end{equation}

\subsubsection{Temperature-Dependent Efficiency Correction}

Utility-scale Li-ion BESS uses an HVAC thermal-management subsystem
(PHVAC) that keeps cells in their optimal window ($15$--$35\,^\circ$C).
Outside this window PHVAC draws a parasitic load that reduces effective
efficiency.
A piecewise-linear modifier $\eta_{{\rm temp},t}\in[0.92,\,1.00]$ is
interpolated from ambient temperature:
%
\begin{equation}
  \eta_{\rm temp}(T) =
  \begin{cases}
    0.94 & T \leq -10\,^\circ\mathrm{C} \\
    0.97 & T = 0\,^\circ\mathrm{C} \\
    1.00 & 15 \leq T \leq 35\,^\circ\mathrm{C} \\
    0.96 & T = 45\,^\circ\mathrm{C} \\
    0.92 & T \geq 60\,^\circ\mathrm{C}
  \end{cases}
  \label{eq:eta_temp}
\end{equation}

At an ERCOT summer peak of $T=40\,^\circ$C ($\eta_{\rm temp}\approx 0.98$),
the effective round-trip efficiency is
$(0.96\times0.98)^2 \approx 88.4\%$.

\subsection{Natural Gas Dispatch Model}

After solar and BESS the residual demand is
%
\begin{equation}
  D_{{\rm res},t}
    = \max\!\bigl(0,\;L_t - P_{{\rm sol},t} - P_{{\rm dis},t}\bigr).
  \label{eq:residual}
\end{equation}

Gas is dispatched as the gap-filler of last resort:
%
\begin{equation}
  P_{{\rm gas},t} = \min\!\bigl(G,\;D_{{\rm res},t}\bigr),
  \label{eq:gas}
\end{equation}

\noindent where $G$ (MW) is installed gas capacity (optimizer decision
variable).
Fuel consumption and CO$_2$ emissions are
%
\begin{align}
  F_t &= P_{{\rm gas},t} \times \mathrm{HR}, \quad [\mathrm{MMBtu}] \label{eq:fuel} \\
  E_{{\rm CO}_2,t} &= F_t \times \varepsilon, \quad [\mathrm{tCO}_2] \label{eq:co2}
\end{align}

\noindent with heat rate $\mathrm{HR} = 9.0$\,MMBtu/MWh and
$\varepsilon = 0.0531$\,tCO$_2$/MMBtu \cite{b4}, giving a full-load
intensity of $0.478$\,tCO$_2$/MWh.

\subsection{Grid-Connected Dispatch (Model B)}

For Model B, residual demand after solar and BESS is split between gas and
grid import based on hourly cost.
The gas short-run marginal cost is
%
\begin{equation}
  c_{\rm src} = \mathrm{HR}\times c_{\rm fuel} + c_{\rm var}
  \quad [\$/\mathrm{MWh}],
  \label{eq:gas_src}
\end{equation}

\noindent where $c_{\rm fuel}$ (\$/MMBtu) is the ISO-specific gas price and
$c_{\rm var} = \$5$/MWh is variable O\&M.
The dispatch each hour is
%
\begin{equation}
  \bigl(P_{{\rm gas},t},\,P_{{\rm grid},t}\bigr) =
  \begin{cases}
    \bigl(0,\;D_{{\rm res},t}\bigr)
      & \mathrm{LMP}_t \leq c_{\rm src} \\[4pt]
    \bigl(\min(G,D_{{\rm res},t}),\;\max(0,D_{{\rm res},t}-G)\bigr)
      & \mathrm{otherwise.}
  \end{cases}
  \label{eq:grid_dispatch}
\end{equation}

% -----------------------------------------------------------------------
\section{Optimization Problem}
\label{sec:opt}

\subsection{Objective Function}

The three optimizer decision variables are $S$ (MW$_{\rm DC}$ solar),
$B$ (MWh battery), and $G$ (MW gas).
The objective is to minimize system levelized cost of energy:
%
\begin{equation}
  \min_{S,\,B,\,G \geq 0}
  \;\mathrm{sLCOE}
  = \frac{C_{\rm sol}(S)+C_{\rm BESS}(B)+C_{\rm gas}(G,E_{\rm gas})}
         {E_{\rm demand}},
  \label{eq:obj}
\end{equation}

\noindent where $E_{\rm demand}$ (MWh/yr) is annual datacenter consumption.

\subsection{Annual Cost Functions}

All costs are annualized using the Capital Recovery Factor
%
\begin{equation}
  \mathrm{CRF}(r,n)
    = \frac{r\,(1+r)^n}{(1+r)^n - 1},
  \label{eq:crf}
\end{equation}

\noindent with discount rate $r=8\%$ and asset-specific life $n$:
%
\begin{align}
  C_{\rm sol}(S)
    &= \bigl[c_{s,c}\,\mathrm{CRF}(r,n_s)
            + c_{s,o}\bigr]\cdot S \cdot 10^3,
  \label{eq:csol} \\[4pt]
  C_{\rm BESS}(B)
    &= \bigl[c_{b,c}\,\mathrm{CRF}(r,n_b)
            + c_{b,{\rm aug}}
            + c_{b,o}/\tau\bigr]\cdot B \cdot 10^3,
  \label{eq:cbess} \\[4pt]
  C_{\rm gas}(G,E_{\rm gas})
    &= \bigl[c_{g,c}\,\mathrm{CRF}(r,n_g)
            + c_{g,f}\bigr]\cdot G \cdot 10^3 \nonumber \\
    &\quad + (c_{g,v} + \mathrm{HR}\cdot c_{\rm fuel})\cdot E_{\rm gas},
  \label{eq:cgas}
\end{align}

\noindent where subscripts $c$, $o$, $f$, $v$ denote CAPEX, fixed O\&M,
fixed O\&M (gas), and variable O\&M, respectively, and $c_{b,{\rm aug}}$
is the annual battery cell-augmentation cost (2.5\%/yr of BESS CAPEX).

For Model B the total cost also includes annualized grid interconnection
infrastructure and the annual grid energy purchase:
%
\begin{equation}
  \begin{split}
    C_{\rm grid}
      &= \bigl[c_{x,c}\,\mathrm{CRF}(r,n_x)+c_{x,o}\bigr]
         \hat{P}_{\rm grid}\cdot 10^3 \\
      &\quad + \sum_{t=1}^{T}\max(0,\mathrm{LMP}_t)\,P_{{\rm grid},t},
  \end{split}
  \label{eq:cgrid}
\end{equation}

\noindent where $\hat{P}_{\rm grid}=\max_t P_{{\rm grid},t}$ (MW) sizes the
interconnection, $c_{x,c}=\$100$/kW, and $c_{x,o}=\$2$/kW/yr.

\subsection{Constraints}

\subsubsection{Reliability Constraint (Model A)}

Because the islanded microgrid must serve load in every hour, the installed
gas capacity must satisfy
%
\begin{equation}
  G \geq G_{\min}(S,B),
  \quad
  G_{\min}(S,B) = \max_{t}\bigl[P_{{\rm gas},t}(S,B,G\!\to\!\infty)\bigr].
  \label{eq:reliability}
\end{equation}

$G_{\min}$ is the peak gas output when gas is unconstrained, equal to the
worst-case hourly residual load the BESS cannot cover.
This surface is precomputed by a two-dimensional grid search before the
economic optimization.

\subsubsection{Renewable Floor Constraint}

A user-specified minimum on-site renewable fraction $\rho_{\min}$ is imposed:
%
\begin{equation}
  \mathcal{R}(S,B,G)
    = \frac{E_{\rm sol,used}+E_{\rm BESS,dis}}{E_{\rm demand}}
    \geq \rho_{\min}.
  \label{eq:ren_floor}
\end{equation}

\noindent Storage dispatch is counted as renewable because the BESS is charged
primarily by solar surplus.

\subsection{Solution Method}

The optimization is solved by exhaustive grid search.
Solar capacity is swept over $S \in \{0,25,\ldots,300\}$\,MW and battery
size over $B \in \{0,40,\ldots,400\}$\,MWh (integer multiples of the 4\,MWh
commercial container), giving 143 candidate $(S,B)$ pairs.
For each pair, $G_{\min}$ is looked up from the precomputed reliability
surface, dispatch is run for all 8,784 hours, sLCOE is computed, and designs
violating \eqref{eq:ren_floor} are excluded.
The grid search avoids local-minimum issues and is fully reproducible.
Total run time is approximately 30\,s on a laptop CPU.

% -----------------------------------------------------------------------
\section{Results}
\label{sec:results}

\subsection{Cross-Market Comparison}

Table~\ref{tab:iso_compare} presents the cost-optimal islanded and
grid-connected designs for all seven ISO/RTO markets at the 50\,MW IT-load
baseline and 30\% renewable floor, ranked by islanded sLCOE.
The Solar and BESS columns show the optimal on-site asset sizes for the
islanded configuration; gas capacity is determined by the reliability
constraint \eqref{eq:reliability} and ranges from 66 to 69\,MW across all
markets.

\begin{table}[htbp]
\caption{All Markets: 50\,MW IT Load, 30\% Renewable Floor (ranked by islanded sLCOE)}
\label{tab:iso_compare}
\begin{center}
\begin{tabular}{|c|c|c|c|c|c|}
\hline
\textbf{ISO} &
\textbf{\begin{tabular}[c]{@{}c@{}}Islanded\\(\$/MWh)\end{tabular}} &
\textbf{\begin{tabular}[c]{@{}c@{}}Grid\\(\$/MWh)\end{tabular}} &
\textbf{\begin{tabular}[c]{@{}c@{}}$\Delta$\\(\$/MWh)\end{tabular}} &
\textbf{\begin{tabular}[c]{@{}c@{}}Solar\\(MW)\end{tabular}} &
\textbf{\begin{tabular}[c]{@{}c@{}}BESS\\(MWh)\end{tabular}} \\
\hline
ERCOT  & 62.14  & 48.97 & 13.17 & 125 & 40 \\
\hline
SPP    & 63.27  & 52.51 & 10.76 & 125 & 40 \\
\hline
MISO   & 71.12  & 63.32 &  7.80 & 150 & 40 \\
\hline
CAISO  & 78.15  & 53.28 & 24.87 & 100 & 40 \\
\hline
PJM    & 78.29  & 60.88 & 17.40 & 175 & 40 \\
\hline
NYISO  & 97.18  & 68.22 & 28.96 & 150 & 80 \\
\hline
ISO-NE & 100.74 & 72.39 & 28.34 & 150 & 80 \\
\hline
\end{tabular}
\end{center}
\end{table}

The grid-connected architecture is cheaper in all seven markets, with the
advantage ($\Delta$) ranging from \$7.80/MWh in MISO to \$28.96/MWh in NYISO.
Three patterns emerge.

First, Sun Belt markets (ERCOT, SPP) achieve the lowest absolute islanded
costs (\$62--\$63/MWh) thanks to their low gas prices (\$2.50--\$2.80/MMBtu),
favorable solar resources, and baseline CAPEX multipliers (1.00$\times$).
However, their wholesale markets are also cheap, so the grid-connected
advantage is moderate (\$10--\$13/MWh).

Second, Northeast markets (NYISO, ISO-NE) have the highest islanded costs
(\$97--\$101/MWh) driven by expensive gas (\$5.50--\$6.00/MMBtu) and high
construction cost multipliers (1.18--1.20$\times$).
These same factors compress the grid-connected savings relative to the
islanded premium, yielding the largest absolute gaps (\$28--\$29/MWh).
They also require 80\,MWh of BESS---double that of Sun Belt ISOs---because
lower solar capacity factors require more intraday shifting to meet the
renewable floor.

Third, CAISO stands out as the largest percentage advantage for the
grid-connected case (32\%): its high gas price inflates the islanded
microgrid while the California wholesale market, despite its volatility,
offers competitively priced import hours.

\subsection{ERCOT Deep Dive}

ERCOT (Dallas-Fort Worth, TX) is selected for a detailed analysis because it
represents the most favorable conditions for islanded operation: lowest gas
price, best solar resource, and lowest CAPEX multiplier.
If grid-connected is still cheaper in ERCOT, it will be cheaper everywhere.

\subsubsection{Effect of the Renewable Constraint}

Table~\ref{tab:ercot} contrasts the unconstrained optimum with the 30\%
renewable floor result for the islanded case.

\begin{table}[htbp]
\caption{ERCOT Deep Dive --- Islanded Optimal Designs}
\label{tab:ercot}
\begin{center}
\begin{tabular}{|l|c|c|}
\hline
\textbf{Metric} & \textbf{No RE floor} & \textbf{30\% RE floor} \\
\hline
Solar $S$ (MW$_{\rm DC}$)  & 0     & 125  \\
\hline
Battery $B$ (MWh)           & 0     & 40   \\
\hline
Gas $G$ (MW)                & 67.6  & 67   \\
\hline
sLCOE (\$/MWh)              & 45.42 & 62.14 \\
\hline
Renewable share (\%)        & 0     & 31   \\
\hline
Annual CO$_2$ (kt)          & 249.7 & 178.5 \\
\hline
\end{tabular}
\end{center}
\end{table}

When no renewable floor is imposed, the unconstrained optimum is pure
gas ($S=0$, $B=0$), driven by the low Texas gas price and the high capital
cost of solar and BESS relative to fuel savings at this price level.
At the 30\% renewable floor the optimal design shifts to $S=125$\,MW DC
with $B=40$\,MWh BESS; the constraint is binding (31\% just meets 30\%),
and sLCOE rises to \$62.14/MWh---a 37\% premium over the unconstrained
gas-only baseline.

\subsubsection{Islanded vs.\ Grid-Connected at 30\% Floor}

At the same 30\% renewable floor the optimal grid-connected design achieves
sLCOE\,=\,\$48.97/MWh, a \$13.17/MWh saving over the islanded alternative.
The grid-connected design benefits from sourcing cheap wholesale imports
during hours when the LMP falls below the on-site gas short-run marginal
cost \eqref{eq:gas_src}, while never being forced to oversize gas capacity
for reliability.
The islanded design must hold 67\,MW of gas firm---sized to the
single worst-case residual demand hour over 8,784 hours---even though that
capacity runs at low utilization most of the year.

% -----------------------------------------------------------------------
\section{Discussion}
\label{sec:discussion}

\subsection{Why the Grid Dominates}

Two structural factors explain the result.
First, the grid acts as an unlimited-capacity backstop: the grid-connected
data center never needs to invest in firm capacity for its own worst-case
load.
The islanded case must size gas to the peak residual demand hour (67.4\,MW
for ERCOT at $S=125$\,MW DC), even when gas runs below 50\% capacity factor
most of the year.
Second, BESS at \$280/kWh provides limited fuel-cost savings relative to its
capital; the optimal islanded designs at the 30\% floor carry only 40\,MWh
of storage in most ISOs (80\,MWh in NYISO and ISO-NE where lower solar
capacity factors require more shifting).

\subsection{Role of the Renewable Constraint}

The renewable floor is binding in all markets---the unconstrained optimum in
every ISO features a renewable share below 30\%.
Any renewable target imposed on an islanded microgrid therefore adds cost
relative to the unconstrained minimum, ranging from roughly \$14--\$28/MWh
in ERCOT and MISO to \$30--\$40/MWh in CAISO, NYISO, and ISO-NE.

\subsection{FERC Order 2222 and Future Economics}

FERC Order No.\ 2222 requires utilities to allow aggregations of distributed
energy resources to participate in wholesale markets.
Full implementation would allow islanded or semi-islanded microgrids to
recover ancillary-service revenues not credited in this model.
BESS CAPEX projected below \$200/kWh by 2030 (NREL Storage Futures) would
materially improve the islanded case economics.

% -----------------------------------------------------------------------
\section{Limitations and Future Work}
\label{sec:limitations}

The current model has several important limitations:

\begin{itemize}
\item \textbf{Deterministic dispatch:} A single 2024 weather/price year is
used.
Stochastic analysis across many historical years would better characterize
tail-risk events such as Winter Storm Uri.

\item \textbf{No unit-commitment:} Gas ramp rates and start-up costs are not
modeled; dispatch below the 30\% minimum stable load is treated as
continuous.

\item \textbf{No demand response:} IT load is non-curtailable.
Hyperscalers may shift batch workloads by minutes to hours in practice.

\item \textbf{No capacity payments:} PJM and ISO-NE operate capacity markets
that could provide revenue credit to firm on-site resources.

\item \textbf{Grid interconnection simplification:} The \$100/kW
interconnection proxy may underestimate upgrade costs in congested zones.
\end{itemize}

Future work includes: (a) multi-year stochastic dispatch; (b) policy
sensitivity over ITC, carbon price, and BESS cost trajectories; (c)
gradient-based continuous optimization for finer asset-size resolution; and
(d) a live dashboard for interactive scenario comparison \cite{b9}.

% -----------------------------------------------------------------------
\section{Conclusions}
\label{sec:conclusions}

This paper presents an open-source hourly simulation framework comparing
islanded BTM and grid-connected energy supply strategies for a 50\,MW data
center across seven U.S.\ ISO/RTO markets.
Key findings are as follows.

\begin{itemize}
\item The grid-connected architecture is cheaper than the islanded microgrid
in all seven markets at the 30\% renewable floor, with advantages ranging
from \$7.80/MWh (MISO) to \$28.96/MWh (NYISO).

\item The renewable floor constraint is binding in all markets: the
unconstrained economic optimum in every region favors gas-dominated or
pure-grid solutions.

\item The primary driver of the islanded microgrid's cost disadvantage is
the requirement to size gas for the worst-case reliability hour, resulting in
low-utilization firm capacity with large fixed charges.

\item Solar without storage is the least-cost path to a 30\% renewable floor
in most markets; battery storage adds cost at current \$280/kWh pricing.

\item Regional gas price and construction cost variation creates a $>$50\%
spread in sLCOE across ISOs for the islanded case.
\end{itemize}

These findings quantify the energy-cost premium a data center operator accepts
in exchange for the resilience that full islanding provides.
As BESS costs continue to fall and FERC Order 2222 matures, the economics of
islanded operation will improve.

% -----------------------------------------------------------------------
\section*{Acknowledgment}

The authors thank the instructors and teaching assistants of ELEN 4510:
Modern \& Clean Energy Systems at Columbia University.
Open data were provided by NREL (PVWatts V8 API and ATB 2025), Open-Meteo
(historical weather), the U.S.\ EIA (STEO and gas prices), and the seven
ISO/RTOs through their public market portals.

% -----------------------------------------------------------------------
\begin{thebibliography}{00}

\bibitem{b1}
International Energy Agency, ``Electricity 2024: Analysis and forecast to
2026,'' IEA Report, Paris, Jan.\ 2024. [Online]. Available:
\url{https://www.iea.org/reports/electricity-2024}

\bibitem{b2}
Schneider Electric, ``Data Center Energy Flexibility,'' White Paper
22511651, 2022. [Online]. Available:
\url{https://www.se.com/us/en/download/document/22511651_EN}

\bibitem{b3}
A.\ P.\ Dobos, ``PVWatts Version 5 Manual,'' NREL Tech.\ Rep.\
NREL/TP-6A20-62641, National Renewable Energy Laboratory, Golden, CO,
2014.

\bibitem{b4}
U.S.\ Environmental Protection Agency, ``AP-42: Compilation of Air
Emissions Factors,'' EPA, Research Triangle Park, NC, 2023. [Online].
Available: \url{https://www.epa.gov/air-emissions-factors-and-quantification/ap-42}

\bibitem{b5}
Z.\ Zippenfenig, ``Open-Meteo.com Weather API,'' GitHub, 2023. [Online].
Available: \url{https://github.com/open-meteo/open-meteo}

\bibitem{b6}
Lazard, ``Lazard's Levelized Cost of Energy+ Analysis---Version 18.0,''
Lazard Ltd., New York, Jun.\ 2024. [Online]. Available:
\url{https://www.lazard.com/research-insights/levelized-cost-of-energyplus}

\bibitem{b7}
National Renewable Energy Laboratory, ``Annual Technology Baseline 2025,''
NREL, Golden, CO, 2025. [Online]. Available: \url{https://atb.nrel.gov}

\bibitem{b8}
U.S.\ Energy Information Administration, ``Short-Term Energy Outlook,''
EIA, Washington, D.C., Apr.\ 2024. [Online]. Available:
\url{https://www.eia.gov/outlooks/steo}

\bibitem{b9}
A.\ Fragkiadaki, D.\ Holland, R.\ Vogeley, and T.\ Mendsaikhan,
``Datacenter Microgrid Simulation,'' GitHub, 2026. [Online]. Available:
\url{https://github.com/fragiadaki/datacenter-microgrid-sim}.
Live dashboard: \url{https://afroditifragiadaki-datacenter-microgrid-sim-dashboard-zfcrhl.streamlit.app/}

\end{thebibliography}

\end{document}
